\begin{frame}
    \frametitle{Problem statement}

    \begin{alertblock}{Initial value problem for Lane-Emden type ODE}
        \begin{equation*}
            \begin{cases}
                \diff[2]{\theta}{\xi}+
                \frac{2}{\xi}
                \diff{\theta}{\xi}+
                N
                \left(\theta\right)=
                g
                \left(\xi\right)
                 & \text { for }0\leq\xi\leq R. \\
                \theta
                \left(0\right)=
                \theta_{0},\quad
                \diff{\theta\left(0\right)}{\xi}=
                \theta^{\prime}_{0}.
                 &
            \end{cases}.
        \end{equation*}
    \end{alertblock}

    \begin{alertblock}{Remarks}
        \begin{description}
            \item[$N\left(\theta\right)=\theta^{m}$]
                  Surge en el modelamiento de la teoría de la
                  estructura estelar.%\cite{umesh}

            \item[$N\left(\theta\right)=\exp\left(\theta\right)$]

                  Aparece en el modelamiento de la difusión de densidad en una esfera de gas isotérmica. %\cite{umesh}

                  Existe una demostración física en el paper de Alzate %\cite{alzate2014}

            \item[$N\left(\theta\right)={\left(y^{2}-C\right)}^{\frac{3}{2}}$]

                  Is turned into the white-dwarf equation, which was introduced in %\cite{chandrasekhar1967} 
                  in his study of the gravitational potential of degenerate stars.
        \end{description}
    \end{alertblock}

    La forma equivalente más usada es:
    \[
        \theta''(\xi) + \frac{2}{\xi}\theta'(\xi) + \theta^n = 0.
    \]

    Con condiciones iniciales:
    \[
        \theta(0) = 1, \qquad \theta'(0)=0.
    \]

    \begin{itemize}
        \item Garantizan regularidad en el origen.
        \item Evitan la singularidad del término \(2/\xi\).
    \end{itemize}

\end{frame}

\begin{frame}{Timeline}
    %\begin{figure}[ht!]
    %   \centering
    %  \includegraphics[width=.7\paperwidth]{timeline}
    %\end{figure}
    % \begin{tikzpicture}[x=1cm, y=1.2cm]
    %     \tikzstyle{etapa}=[circle, draw, fill=blue!20, minimum size=0.7cm]
    %     \tikzstyle{texto}=[align=left]
    %     \node[etapa] (p1) at (0,0) {};
    %     \node[texto] at (2,0) {\textbf{1900}\\Primeros modelos};
    %     \node[etapa] (p2) at (0,-1.5) {};
    %     \node[texto] at (2,-1.5) {\textbf{1930}\\Chandrasekhar};
    %     \node[etapa] (p3) at (0,-3) {};
    %     \node[texto] at (2,-3) {\textbf{1967}\\Stellar Structure};
    %     \node[etapa] (p4) at (0,-4.5) {};
    %     \node[texto] at (2,-4.5) {\textbf{2025}\\Modelos actuales};
    %     \draw[thick] (p1) -- (p4);
    % \end{tikzpicture}
    % La principal dificultad al resolver ecuaciones diferenciales singulares es el
    % comportamiento singular que ocurre en \(r = 0\). Sin embargo, se han utilizado
    % muchos métodos para resolver este tipo de ecuaciones diferenciales singulares.
    % Wazwaz (2005)
    % Khan (2012) aplicaron el método de transformación diferencial (DTM)
    % Chowdhury y Hashim (2009) utilizaron el método de perturbación por homotopía (HPM).
    % Baranwal \textit{et al}. (2012) aplicaron un enfoque híbrido
    % Parand \textit{et al}. (2009) utilizaron una técnica seudoespectral basada en las funciones de Legendre
    % racionales y la integración de Gauss--Radau para la solución de la ecuación de
    % Lane--Emden de tipo (1).

    % Saeed Mutaish y Hasan (2018), y Hasan y Saeed Mutaish (2018)
    % resolvieron ecuaciones de tipo Lane--Emden mediante un método de descomposición
    % de Adomian ajustado y también describieron la formación de estas ecuaciones.
    % %\cite{umesh}
\end{frame}